\section{Absolute Scale}
\label{sec:scale}

Pure geometry supplies one length $s$ and pure numbers. Stress, energy and mass
require a dimensionful bridge. We adopt $\hbar c$ as that external constant
\assumed:
\[
|\sigma_{\mathrm{res}}|=C\,\frac{\hbar c}{s^2},
\]
with $C$ a pure geometric prefactor of order $10$--$50$.

Reading residual-stress energy of the membrane volume as rest energy
\interpretive\ gives
\[
m\sim\frac{\hbar}{c\,s}.
\]
The fundamental length is fixed at the proton charge radius
\interpretive:
\[
s\approx 0.84\,\mathrm{fm}.
\]
With $\hbar c\approx 197.3\,\mathrm{MeV\cdot fm}$ one obtains
\[
\frac{\hbar c}{s^2}\approx 280\,\mathrm{MeV\cdot fm^{-1}},\qquad
m\sim 235\,\mathrm{MeV}/c^2
\]
(order of light-hadron masses), and
\[
\ell\approx 0.27\,\mathrm{fm}.
\]

\begin{remark}
Identifying $s$ with the Planck length would close the framework onto Planck units
but place the particle at the Planck mass. The hadronic choice is preferred if the
layered cell is to describe ordinary matter.
\end{remark}
